\section{A uniform projected Hermitian--Einstein equation}
\label{sec:projected-he}

We now put the effective holomorphic bundle family in controlled
coordinates and solve the infinite-dimensional part of the HE
equation. The estimates use the section Laplacian and first
derivatives. They do not require a uniform bound for the full
curvature of the degenerating background.

\subsection{The section inverse in a graph norm}

Write $P_p=\dbar_0^*\dbar_0=\tfrac12D_0^*D_0$ on complex
endomorphism sections, and let $\widehat\Pi_p$ be projection
onto its parallel kernel $\widehat{\mathfrak h}
=\mathfrak{gl}_2\oplus\mathfrak{gl}_2$.
The unrestricted graph space is
\[
 \mathcal G_p=\{s\in W^{1,\infty}:P_ps\in L^\infty\},
 \qquad
 \|s\|_{\mathcal G}=\|s\|_\infty+\|D_0s\|_\infty+\|P_ps\|_\infty,
\]
where the equation is distributional. On trace-free Hermitian
sections write $\Pi_p$ for projection onto
$\mathfrak h$ in \eqref{eq:seven-space}, and put
\begin{equation}\label{eq:real-graph-space}
 \mathcal X_p^0=\{s\in\mathcal G_p:s=s^{*h_0},\ \tr s=0,
                                          \ \Pi_ps=0\}.
\end{equation}

\begin{lemma}\label{lem:graph-inverse}
The following bounds hold uniformly on the full smooth period
base near zero:
\begin{equation}\label{eq:graph-inverse}
 \|s\|_\infty+\|D_0s\|_\infty
                        \le C\|P_ps\|_\infty,
                  \qquad\widehat\Pi_ps=0.
\end{equation}
Consequently $P_p^{-1}$ maps the $L^\infty$ complement uniformly
into its graph domain, and in particular maps the real
trace-free complement into $\mathcal X_p^0$.
The spaces are Banach spaces; $\mathcal G_p$ is a Banach algebra
with a uniform product constant. Moreover
$\|\eta_p^\pm\|_\infty\le C$, and harmonic projection in
endomorphism degree one is uniformly bounded on $L^\infty$.
\end{lemma}

\begin{proof}
For scalar heat flow $H_t=e^{-t\Delta}$, Bochner's formula
and nonnegative Ricci curvature give
\begin{equation}\label{eq:reverse-poincare}
 2t|\dd H_tv|^2\le H_t(v^2)-(H_tv)^2
                                        \le\|v\|_\infty^2.
\end{equation}
Indeed differentiation of $H_{t-a}((H_av)^2)$ yields
\[
 H_t(v^2)-(H_tv)^2
       =2\int_0^t H_{t-a}(|\dd H_av|^2)\,\dd a.
\]
The Bochner gradient contraction bounds each integrand below
by $|\dd H_tv|^2$. This also proves the inequality in the
normalization of \cite[Section 2, Proposition 3]{Ledoux}.
Taking real and imaginary parts and integrating
$w=H_Tw+\int_0^T H_t\Delta w\,\dd t$ gives
\begin{equation}\label{eq:scalar-gradient-estimate}
 \|\dd w\|_\infty\le CT^{-1/2}\|w\|_\infty
                              +C\sqrt T\|\Delta w\|_\infty.
\end{equation}
The $t^{-1/2}$ derivative bound is integrable at zero.
Applied to the bounded line potential of
Proposition~\ref{prop:line-he}, this proves
$\|\dd u_p\|_\infty\le C$.

Choose a fixed finite ambient holomorphic frame cover for
$\OO(1,-1)$ and a subordinate partition with supports
compactly contained in the frame charts. The restrictions
of its cutoffs have uniformly bounded gradients and scalar
Laplacians. The first bound follows from
\eqref{eq:diameter-ambient}; the second follows from the
Kähler trace of the fixed ambient complex Hessian.
In the corresponding unitary line frame,
\[
 a_p=\tfrac12(\partial-\dbar)\log h_p=\mathrm i\vartheta_p
\]
is uniformly bounded. It also satisfies $\delta\vartheta_p=0$:
on a Kähler chart $\delta\dd^c v=0$ because the parallel
complex structure contracts the symmetric real Hessian
antisymmetrically. The same statements hold for the
$\pm2$ tensor powers. For a scalar coefficient $v$ of a
section, their rough operator is
\begin{equation}\label{eq:local-rough-operator}
 D_0^*D_0v=\Delta v-2\mathrm i\langle\vartheta_p,\dd v\rangle
                +(\mathrm i\delta\vartheta_p+|\vartheta_p|^2)v.
\end{equation}
Its first-order and zeroth-order coefficients are therefore
uniformly bounded, with the divergence term zero.

Parabolic Kato and the maximum principle give
\[
 |e^{-tP_p}f|\le H_{t/2}|f|.
\]
This follows by applying the scalar inequality to a
regularized norm for the equation
$\partial_ts=-\tfrac12D_0^*D_0s$; it uses no differentiated
bundle domination assertion. At any fixed positive time
the scalar heat kernel is uniformly bounded: apply the
Li--Yau Harnack inequality to compare $H_T(x,y)$ with
$H_{2T}(x,z)$ and integrate over $z$, using fixed volume
and bounded diameter. Thus $e^{-P_p}$ maps $L^2$ uniformly
to $L^\infty$. On $\widehat{\mathfrak h}^{\perp}$ the
section gap gives, for $t\ge1$,
\[
 \|e^{-tP_p}f\|_\infty
          \le Ce^{-c(t-1)}\|f\|_2
          \le C'e^{-c(t-1)}\|f\|_\infty.
\]
Use contraction up to time one and integrate the semigroup
to obtain $\|P_p^{-1}f\|_\infty\le C\|f\|_\infty$.

For gradient control, localize each scalar line coefficient
$v_j$ by its ambient cutoff $\chi_j$ and extend
$w_j=\chi_jv_j$ by zero to the fibre. Equation
\eqref{eq:local-rough-operator} gives
\[
 \|\Delta w_j\|_\infty\le C\bigl(\|P_ps\|_\infty
                  +\|D_0s\|_\infty+\|s\|_\infty\bigr).
\]
Apply \eqref{eq:scalar-gradient-estimate} and sum over the
fixed number of charts and matrix entries. Since the
cutoffs sum to one and the connection forms are bounded,
\[
 \begin{split}
 \|D_0s\|_\infty\le{}&C(T^{-1/2}+1)\|s\|_\infty\\
 &+C\sqrt T\bigl(\|P_ps\|_\infty
                        +\|D_0s\|_\infty+\|s\|_\infty\bigr).
 \end{split}
\]
Choose one small $T$ to absorb the gradient on the right.
Combined with the sup-norm inverse, this is
\eqref{eq:graph-inverse}.

Heat-smoothing a bounded right side preserves the
orthogonal complement and its sup-norm bound. On each
fixed smooth fibre, the inverses have weak-star limits in
$W^{1,\infty}$ which solve the equation distributionally;
uniqueness identifies them with the spectral inverse.
The estimates pass to the limit. The graph spaces are
complete because a Cauchy sequence converges in
$W^{1,\infty}$ and its $P_p$ images converge in $L^\infty$;
distributional continuity identifies the two limits.
The unconstrained kernel is
$\widehat{\mathfrak h}$, and its intersection with the
orthogonal graph domain is zero. On the full graph space,
\begin{equation}\label{eq:graph-product}
 P_p(uv)=(P_pu)v+u(P_pv)
                   -\sum_\ell(D_{0,e_\ell}u)(D_{0,e_\ell}v)
\end{equation}
proves the uniform algebra bound. The identity belongs to
this unrestricted space. Product and chain rules on the
weak domain follow from fixed-fibre $W^{2,q}$ regularity
for every finite $q$; no uniform $W^{2,q}$ bound is needed.

Finally the right sides in
\eqref{eq:harmonic-section-projection} are bounded and
supported on a fixed exterior. The gradient inverse gives
$\|\eta_p^\pm\|_\infty\le C$. The eight harmonic matrix
basis vectors have masses bounded below by
\eqref{eq:extension-coefficients}. Their rank-one
projection formulas, with fixed volume and bounded
representatives, give the $L^\infty$ projection bound.
\end{proof}

\subsection{Controlled coordinates on the actual bundle family}

Let $K=U(2)\times U(2)$ act by
\eqref{eq:two-arrow-action}. The equivariant germ in
Proposition~\ref{prop:effective-bundles} has a representative
on a common $K$-invariant matrix ball $|z|<a$. Indeed the
étale construction has an actual group action before local
inversion; uniqueness of its branch makes the local inverse
and its $K$-translates agree, and compactness of $K$ permits
one common domain. This retains the independent full period
ball after shrinking it.

We construct determinant-preserving smooth identifications
of these bundles with the split bundle, holomorphic in $z$,
which are holomorphic in the space variable on smaller fixed
node neighbourhoods. The total conifold space is smooth,
and the actual bundle is locally free across its central
nodes. On a finite cover of a neighbourhood of the compact
central fibre, choose local holomorphic bundle frames.
On each node neighbourhood these give an identification
with the split bundle, equal to the identity at $z=0$.
Average it over $K$. The average remains holomorphic in
space and parameter and is the identity at zero, hence
invertible on a smaller common ball. It is $K$-equivariant.

Use these frames to define connections only in the matrix
parameter directions. Patch them by a partition depending
on the total-space point and not on $z$, choosing the node
connection exclusively on each smaller core. Average this
parameter connection over $K$ as well. Its coefficients
remain holomorphic in $z$. Parallel transport along
$t\mapsto tz$, $0\le t\le1$, gives an identification $T_z$
which is smooth in space, holomorphic in $z$, and equivariant.
The holomorphic parameter dependence follows directly from
the iterated-integral solution of its linear ODE.
On the cores it agrees with the chosen holomorphic
identification. The determinant trivialization can also
be chosen equivariantly, with its fixed determinant
character at zero, by averaging and shrinking.
If $\det T_z$ differs from it by $\delta(x,z)$, then
$\delta(x,0)=1$. Replacing $T_z$ by $T_z\delta^{-1/4}$,
with the branch near one, preserves these properties and
fixes the determinant exactly.

Consequently the actual holomorphic structures become
$\dbar_0+\beta_p(z)$, where $\tr\beta_p(z)=0$,
$\beta_p(0)=0$, and $\beta_p(z)$ vanishes on the fixed
cores. Compactness of the exterior construction gives
uniform spatial and matrix-parameter bounds of every
fixed order. In particular
\begin{equation}\label{eq:initial-beta-bounds}
 \|\beta_p(z)\|_\infty+
             \|\dbar_0^*\beta_p(z)\|_\infty\le C|z|.
\end{equation}
Normalize the two scalar Kodaira--Spencer frames to the
Koszul classes $e_\pm$. A required change is a nonvanishing
holomorphic scalar function of $p$, absorbed in the arrow
coordinate on a smaller common ball. It leaves the base
coordinates and the two-arrow action unchanged.
Write $\beta_{1,p}(z)=d_z\beta_p(0)z$.
Equivariance under $(tI_2,I_2)$ forces its diagonal
blocks to vanish and its two off-diagonal blocks to have
the corresponding arrow weights.

Define the linear endomorphism and complex gauge
\[
 f_p(z)=P_p^{-1}\dbar_0^*\beta_{1,p}(z),
                  \qquad g_p(z)=e^{-f_p(z)}.
\]
The inverse is on the off-diagonal blocks, whose kernel
is zero. Lemma~\ref{lem:graph-inverse} gives
$\|f_p(z)\|_{\mathcal G}\le C|z|$.
It is trace-free, so the gauge has determinant one.
The transformed coefficient
\begin{equation}\label{eq:harmonic-gauge}
 b_p(z)=g_p(z)^{-1}\beta_p(z)g_p(z)
                         +g_p(z)^{-1}\dbar_0g_p(z)
\end{equation}
is integrable, trace-free and equivariant. Its linear
part is $\beta_{1,p}-\dbar_0f_p$, which is closed,
coclosed and in the prescribed class, hence equals
$\eta_p(z)$ in \eqref{eq:harmonic-arrow-form}.

The graph algebra \eqref{eq:graph-product} controls the
exponential remainder, including its $P_p$ image:
\[
 \|g_p-I+f_p\|_{\mathcal G}
              +\|g_p^{-1}-I-f_p\|_{\mathcal G}\le C|z|^2.
\]
For multiplication of a $(0,1)$-form by an endomorphism,
the product rule for $\dbar_0^*$ has error bounded by
the product of the endomorphism's first derivative and
the form itself. Apply this on both sides of
\eqref{eq:harmonic-gauge}. The linear terms are
$\dbar_0^*\beta_{1,p}$ and $-P_pf_p$ and cancel exactly.
We have therefore proved
\begin{equation}\label{eq:controlled-harmonic-gauge}
 \begin{gathered}
 b_p(z)=\eta_p(z)+N_p(z),\qquad \|b_p(z)\|_\infty\le C|z|,\\
 \|N_p(z)\|_\infty+\|\dbar_0^*b_p(z)\|_\infty\le C|z|^2,
 \qquad\|d_zN_p(z)\|_\infty\le C|z|.
 \end{gathered}
\end{equation}
Holomorphic transport and the uniformly convergent
exponential series give bounds for every fixed matrix
derivative on a smaller common ball. The contracted
derivative estimate uses $P_pf_p$, not a full Hessian
bound for $f_p$.

\subsection{The projected equation}

With $b=b_p(z)$, the Chern connection for $h_0$ has
$D_b''=\dbar_0+b$ and $D_b'=D_0'-b^*$.
On endomorphisms these additional terms act by commutators.
For $s\in\mathcal X_p^0$ define
\begin{equation}\label{eq:actual-he-map}
 \mathcal F_p(s,z)=e^{s/2}
       (\mathrm i\Lambda F_{h_0e^s,b})e^{-s/2}.
\end{equation}
The conjugation makes the output $h_0$-Hermitian, and
its trace is zero because the determinant is fixed.

\begin{proposition}\label{prop:projected-he}
There are uniform small radii such that, for every $p$
and $z$ in their domains, a unique small
$s_p(z)\in\mathcal X_p^0$ solves
\begin{equation}\label{eq:projected-equation}
             (I-\Pi_p)\mathcal F_p(s_p(z),z)=0.
\end{equation}
The solution is smooth on each fibre, $K$-equivariant,
and smooth in the real matrix parameters, with
\begin{equation}\label{eq:he-complement-size}
 \|s_p(z)\|_{\mathcal G}\le C|z|^2,
 \qquad\|d_zs_p(z)\|_{\mathcal G}\le C|z|.
\end{equation}
Every fixed higher matrix derivative has a uniform bound
on a smaller common ball. The remaining actual equation is
\begin{equation}\label{eq:finite-actual-moment}
 \Psi_p(z):=\Pi_p\mathcal F_p(s_p(z),z)=0,
           \qquad\Psi_p(z)\in\mathfrak h,
           \quad\|\Psi_p(z)\|\le C|z|^2.
\end{equation}
The metric $h_0e^{s_p(z)}$ is HE precisely on this zero set.
\end{proposition}

\begin{proof}
Put $M_b=\mathrm i\Lambda F_{h_0,b}$.
Before contraction its curvature difference is
\[
 D_0'b-\dbar_0b^*-b\wedge b^*-b^*\wedge b.
\]
The contracted linear terms are controlled by
$\dbar_0^*b$ and its adjoint, by the Kähler identities.
Thus \eqref{eq:controlled-harmonic-gauge} gives
$\|M_b\|_\infty\le C|z|^2$.
The exact finite change of Chern curvature is
\begin{equation}\label{eq:finite-chern-change}
 \mathrm i\Lambda F_{h_0e^s,b}
       =M_b+\mathrm i\Lambda\dbar_b(e^{-s}D_b'e^s).
\end{equation}
This identity, with the displayed matrix order, is
\cite[(2.10)]{CLZZ}; its infinitesimal version is
\cite[(6.3.16)--(6.3.17)]{Kobayashi}.
At $(s,z)=(0,0)$ the derivative in $s$ is
\[
 D_s\mathcal F_p(0,0)s
       =\mathrm i\Lambda\dbar_0D_0's=P_ps.
\]
For scalars the normalization is
$\delta(\mathrm i\Lambda F)=\tfrac12\Delta s$ for
$h\mapsto he^s$.

For a general endomorphism $v$, expand
\[
 \begin{split}
 \dbar_bD_b'v={}&\dbar_0D_0'v-[\dbar_0b^*,v]
             +[b^*,\dbar_0v]\\
             &+[b,D_0'v]-[b,[b^*,v]],
 \end{split}
\]
with graded brackets for one-forms. After contraction,
\begin{equation}\label{eq:contracted-perturbation}
 \begin{split}
 \|\mathrm i\Lambda\dbar_bD_b'v-P_pv\|_\infty
 \le C\bigl(&\|b\|_\infty\|D_0v\|_\infty\\
 &+(\|\dbar_0^*b\|_\infty+\|b\|_\infty^2)\|v\|_\infty\bigr).
 \end{split}
\end{equation}
No uncontracted derivative of $b$ occurs.
In \eqref{eq:finite-chern-change}, expand
\[
 \dbar_b(e^{-s}D_b'e^s)
   =(\dbar_be^{-s})\wedge D_b'e^s
                            +e^{-s}\dbar_bD_b'e^s.
\]
The first term uses only first derivatives. Control the
second by \eqref{eq:contracted-perturbation} and the graph
product rule. The exponential series and conjugation in
\eqref{eq:actual-he-map} are uniformly smooth maps in
these norms. Their weak-domain formulas are legitimate
by fixed-fibre $W^{2,q}$ regularity, without any assertion
of density in the graph norm.

Writing $\mathcal F_p(s,z)=P_ps+\mathcal R_p(s,z)$ gives
\begin{equation}\label{eq:he-remainder}
 \begin{aligned}
 \|\mathcal R_p(0,z)\|_\infty&\le C|z|^2,\\
 \|\mathcal R_p(s,z)-\mathcal R_p(t,z)\|_\infty
  &\le C(|z|+\|s\|_{\mathcal G}+\|t\|_{\mathcal G})
                                      \|s-t\|_{\mathcal G}.
 \end{aligned}
\end{equation}
To check the vanishing factor in the second line, every
term beyond $P_ps$ contains $b=O(|z|)$, an exponential
remainder factor $s$ or $t$, or a first derivative of
one of them. The contracted derivative of $b$ is
$O(|z|^2)$. The product series and their differentiated
series prove the estimate and bounded higher derivatives.

Projection $\Pi_p$ is uniformly bounded on $L^\infty$,
since it is averaging against parallel block matrices
with the fixed mass \eqref{eq:split-mass}.
Let $J_p$ be the real complementary inverse from
Lemma~\ref{lem:graph-inverse}. Equation
\eqref{eq:projected-equation} is equivalent to
\[
       s=-J_p(I-\Pi_p)\mathcal R_p(s,z).
\]
Choose a fixed small graph ball and then a common matrix
ball so that its Lipschitz constant is less than $1/2$.
The quadratic forcing makes the ball of radius
$C_0|z|^2$ invariant. Contraction gives the solution and
its quadratic bound. Every solution in the original
small graph ball satisfies the same bound after
absorption, and hence is this fixed point.
Differentiating the uniformly smooth fixed-point equation
gives the higher matrix bounds; $s_p(0)=d_zs_p(0)=0$
gives the derivative estimate in
\eqref{eq:he-complement-size}. All operations commute
with $K$, so uniqueness proves equivariance.

For each fixed smooth fibre the graph solution is in
$W^{2,q}$ for all finite $q$, hence in $C^{1,\alpha}$.
The coefficient of its contracted second derivative in
\eqref{eq:actual-he-map} is
\[
 v\longmapsto e^{s/2}
       \left(\int_0^1 e^{-ts}v e^{ts}\,\dd t\right)e^{-s/2}.
\]
It is the identity at zero and has a smooth pointwise
inverse for small $s$. The equation therefore expresses
$P_ps$ as a $C^{0,\alpha}$ function of first derivatives
and smooth fixed-fibre data. Schauder regularity and
iteration make $s$ smooth. This bootstrap uses no uniform
higher unweighted spatial estimate. Finally
$\Pi_pP_ps=0$ and \eqref{eq:he-remainder} give the size
of $\Psi_p$. The projected equation makes the full
transported curvature moment equal to $\Psi_p$, proving
the final equivalence.
\end{proof}

The finite equation \eqref{eq:finite-actual-moment} will
be solved below. At this stage the construction is an
actual complementary metric correction for every small
integrable matrix parameter. Its constants control matrix
derivatives and the stated spatial norms throughout
$\mathcal B^\circ$; they make no claim about derivatives
in the degenerating base parameter.
